% ============================================================================
%  MIT 805 (2026) - Part 1, Section 5: The 7 Vs of Big Data
%  Dataset: NYC TLC High-Volume For-Hire Vehicle (HVFHV) Trip Record Data
%
%  Preamble needs: \usepackage{booktabs}
%                  \usepackage{microtype}   % recommended, absorbs long \texttt{}
% ============================================================================

\section{The 7 Vs of Big Data}
\label{sec:7vs}

The TLC High-Volume For-Hire Vehicle (HVFHV) records are a seven-year,
individually resolved account of app-dispatched urban mobility.

\begin{table}[htbp]
\centering\footnotesize
\caption{The 7 Vs applied to the NYC TLC HVFHV Trip Record Data}
\label{tab:7vs}
\begin{tabular}{@{}p{1.5cm}p{7.3cm}p{7.3cm}@{}}
\toprule
\textbf{V} & \textbf{Evidence in our dataset} & \textbf{Why it matters here} \\
\midrule

Volume &
Raw 37.58\,GB across 89 monthly Parquet files (Feb 2019--Jun 2026); June 2026
alone holds 20.78 million trips. &
A single month is the practical ceiling for in-memory \texttt{pandas} work;
the $\geq$12\,GB working and $\geq$3\,GB processing sets exceed it, so Part~2
requires partitioned, distributed (PySpark) execution. \\
\addlinespace

Velocity &
${\sim}692{,}500$ trips/day in June 2026 ($\approx$\,\textbf{8 dispatches/s});
each record carries four event timestamps (request, on-scene, pickup,
dropoff). &
Fast at source but published monthly: the data suits month-partitioned batch
ingestion and retrospective event-time analysis, not live dispatch. \\
\addlinespace

Variety &
25 variables of mixed kind: timestamps; continuous measures
(\texttt{trip\_miles}, \texttt{trip\_time}); money fields
(\texttt{base\_passenger\_fare}, \texttt{driver\_pay}, \texttt{tips}); Y/N
flags; license codes (Uber HV0003, Lyft HV0005); zone IDs
(\texttt{PULocationID}/\texttt{DOLocationID}). &
Heterogeneity drives pre-processing cost: the money fields are not comparable
(fare excludes tolls/tips/taxes; \texttt{driver\_pay} is net of commission),
flags and codes are categories, and zone IDs are meaningless until joined to
the 265-zone lookup and shapefile. \\
\addlinespace

Veracity &
TLC disclaims accuracy. June 2026 checks: 4{,}844 zero-distance, 6{,}750
negative-fare and 2 non-positive-duration trips; 14{,}421 records
(\textbf{0.07\%}) excluded, leaving 20{,}761{,}447. &
The defects are systematic, not noise: negative fares (refunds/voids) flip the
sign of fare aggregates, and zero-distance trips corrupt speed statistics, so
records are filtered under a documented rule. \\
\addlinespace

Variability &
\emph{Schema:} fare/\texttt{driver\_pay} columns added May 2022;
\texttt{cbd\_congestion\_fee} only from Jan 2025. \emph{Behaviour:} mean
distance 7.13\,mi at 05:00 vs 4.70\,mi at 15:00 (\textbf{+51.7\%}); Saturday
+19.3\% trips over Wednesday. &
Neither schema nor behaviour stays constant: analyses are only valid where the
columns exist, and one 89-month average would mix COVID-19, pre- and
post-congestion-pricing periods, so results must be computed per period. \\
\addlinespace

Visualiz\-ation &
Borough pickups (Manhattan 7.14m, Brooklyn 5.72m, Queens 4.77m, Bronx 2.79m,
Staten Island 0.35m); hourly distance/duration profiles; LaGuardia 421{,}225
vs JFK 356{,}638; Uber 73.4\%\,/\,Lyft 26.6\%; $265{\times}265$ OD matrix. &
No reader can inspect 20.78m rows, so findings are delivered as graphics; the
zero-distance and negative-fare defects were first visible as artefacts in the
plotted distributions. \\
\addlinespace

Value &
\texttt{driver\_pay}, \texttt{tips} and wait times support effective-earnings
analysis; request-vs-match flags measure pooling and unmet
wheelchair-accessible demand. \textbf{65.6\% of pickups occur outside
Manhattan.} &
Demand is strongest where the subway is weakest, informing transit-gap,
curb-space and equity decisions; earnings fields quantify driver welfare; and
the Jan 2025 congestion-fee column gives a natural before/after policy
window. \\
\bottomrule
\end{tabular}
\end{table}
